\begin{lemma}
In signature form, suppose a rooted tree has at most $\Delta$ children at
each vertex, at most $h$ marked children at each vertex, and at most $w$
unmarked vertices on every root-to-vertex path.  If $k\geq w$ and
$h\leq\Delta$, let $\operatorname{count}(z)$ be the number of level-$k$
paths with binary signature $z$.  Thus
$\operatorname{count}(z)\leq\Delta^j h^{k-j}$ for a signature with $j$
unmarked positions, and it is zero for $j>w$.  Then the total number of
root-to-level-$k$ paths is at most
\[
 2^k\Delta^w h^{k-w}.
\]
Equivalently, after recording each path by its binary marked/unmarked
signature, the number of paths with a signature containing $j$ unmarked
positions is at most $\Delta^j h^{k-j}$, and only signatures with
$j\leq w$ can occur.
\end{lemma}
\begin{proof}
Fix a binary signature $z$ and let $j$ be its number of unmarked positions.
Read a path from the root.  At an unmarked position there are at most
$\Delta$ choices, and at a marked position at most $h$ choices.  Induction
on the length of the path therefore bounds the number of paths with this
signature by $\Delta^j h^{k-j}$.  The path hypothesis makes this number zero
when $j>w$.  Summing over the at most $2^k$ binary signatures and using
$h\leq\Delta$ gives, for every $j\leq w$,
$\Delta^j h^{k-j}\leq\Delta^w h^{k-w}$; hence the total is at most
$2^k\Delta^w h^{k-w}$.
\end{proof}
